% 【拆分生成】 answers.tex —— 2022 年·新高考Ⅰ卷 · 参考答案与详细解析片段
% 【拆分生成】 由 tools/split_project.py 自合并书 answers/2022_xgk1_sol.tex 包裹 answer 环境生成，勿手改
% 由 tools/md2tex.py 自 2022_数学真题_含解析.md 生成，勿手改（改源 MD 后重跑）


\begin{gknotice}
\textbf{解析归属}：本文档全部解析由本库编写者根据试题独立推导，\textbf{非官方标准答案}。推导过程完整给出以便核验；若与官方发布版本不一致，以官方为准。
\end{gknotice}

\par\medskip

\gkpart{一、单选题（每小题 5 分，共 40 分）}

\begin{answer}
\gksolhead{1}{}

\textbf{题目}　若集合 $M=\left\{x\mid\sqrt x<4\right\}$，$N=\{x\mid3x\ge1\}$，则 $M\cap N=$（　　）

\textbf{解析}

\textbf{求 $M$（含一个极易遗漏的隐含条件）}：

$\sqrt x$ 有意义要求 $x\ge0$。在此前提下，$\sqrt x<4$ 两边平方（均非负，不等号方向不变）：
\[x<16\]

故
\[M=\{x\mid0\le x<16\}\]

\textbf{求 $N$}：
\[3x\ge1\ \Longrightarrow\ x\ge\frac13\]

\textbf{取交集}：
\[M\cap N=\left\{x\ \middle|\ \frac13\le x<16\right\}\]

\gkanswer{答案：D}

\begin{gktip}
\textbf{易错点}：忘记 $\sqrt x$ 的定义域 $x\ge0$，得到 $M=\{x\mid x<16\}$，交集变成 $\left\{x\mid\dfrac13\le x<16\right\}$——本题恰好答案相同（因为 $\dfrac13>0$ 已经保证非负），但在其它题目中这类疏漏往往致命。\textbf{看到根号，先写定义域}应成为条件反射。
\end{gktip}

\par\medskip

\end{answer}
\begin{answer}
\gksolhead{2}{}

\textbf{题目}　若 $\mathrm i(1-z)=1$，则 $z+\bar z=$（　　）

\textbf{解析}

\textbf{第一步，解出 $z$。}
\[1-z=\frac1{\mathrm i}\]

利用 $\dfrac1{\mathrm i}=-\mathrm i$：
\[1-z=-\mathrm i\ \Longrightarrow\ z=1+\mathrm i\]

\textbf{第二步，求 $z+\bar z$。}

$\bar z=1-\mathrm i$，故
\[z+\bar z=(1+\mathrm i)+(1-\mathrm i)=2\]

\gkanswer{答案：D}

\begin{gktip}
\textbf{通用结论}：$z+\bar z=2\operatorname{Re}(z)$。本题中 $\operatorname{Re}(z)=1$，直接得 $2$，连 $\bar z$ 都不必写出。这个结论在选填题中能省下不少时间。
\end{gktip}

\par\medskip

\end{answer}
\begin{answer}
\gksolhead{3}{}

\textbf{题目}　在 $\triangle ABC$ 中，点 $D$ 在边 $AB$ 上，$BD=2DA$。记 $\overrightarrow{CA}=\boldsymbol m$，$\overrightarrow{CD}=\boldsymbol n$，则 $\overrightarrow{CB}=$（　　）

\textbf{解析}

\textbf{第一步，弄清 $D$ 的位置（本题唯一的陷阱）。}

$BD=2DA$ 意味着 $D$ 到 $B$ 的距离是到 $A$ 的两倍，即 $D$ \textbf{靠近 $A$}，且
\[AD:DB=1:2\ \Longrightarrow\ \overrightarrow{AD}=\frac13\overrightarrow{AB}\]

\textbf{第二步，用向量的三角形法则建立关系。}
\[\overrightarrow{CD}=\overrightarrow{CA}+\overrightarrow{AD}=\boldsymbol m+\frac13\overrightarrow{AB}\]

而
\[\overrightarrow{AB}=\overrightarrow{CB}-\overrightarrow{CA}=\overrightarrow{CB}-\boldsymbol m\]

代入：
\[\boldsymbol n=\boldsymbol m+\frac13\left(\overrightarrow{CB}-\boldsymbol m\right)=\frac23\boldsymbol m+\frac13\overrightarrow{CB}\]

\textbf{第三步，解出 $\overrightarrow{CB}$。}
\[\frac13\overrightarrow{CB}=\boldsymbol n-\frac23\boldsymbol m\ \Longrightarrow\ \overrightarrow{CB}=3\boldsymbol n-2\boldsymbol m=-2\boldsymbol m+3\boldsymbol n\]

\gkanswer{答案：B}

\begin{gktip}
\textbf{验证技巧}：把结果代回检验——若 $\overrightarrow{CB}=3\boldsymbol n-2\boldsymbol m$，则 $\overrightarrow{AB}=\overrightarrow{CB}-\boldsymbol m=3\boldsymbol n-3\boldsymbol m$，$\dfrac13\overrightarrow{AB}=\boldsymbol n-\boldsymbol m=\overrightarrow{AD}$  与 $\overrightarrow{CD}=\overrightarrow{CA}+\overrightarrow{AD}=\boldsymbol m+(\boldsymbol n-\boldsymbol m)=\boldsymbol n$ 吻合。
\end{gktip}

\par\medskip

\end{answer}
\begin{answer}
\gksolhead{4}{}

\textbf{题目}　水库在水位 $148.5\ \mathrm m$ 时水面面积 $140.0\ \mathrm{km}^{2}$，水位 $157.5\ \mathrm m$ 时水面面积 $180.0\ \mathrm{km}^{2}$。将该形状视为棱台，求水量增加值。（$\sqrt7\approx2.65$）

\textbf{解析}

\textbf{第一步，确定棱台的三个要素。}
\[h=157.5-148.5=9\ (\mathrm m)\]
\[S_1=140.0\ \mathrm{km}^{2},\qquad S_2=180.0\ \mathrm{km}^{2}\]

\textbf{第二步，计算几何平均项。}
\[\sqrt{S_1S_2}=\sqrt{140\times180}=\sqrt{25200}\]

分解：$25200=3600\times7$，故
\[\sqrt{25200}=60\sqrt7\approx60\times2.65=159\ (\mathrm{km}^{2})\]

（题目给 $\sqrt7\approx2.65$ 正是为这一步准备的。）

\textbf{第三步，代入棱台体积公式。}
\[V=\frac h3\left(S_1+S_2+\sqrt{S_1S_2}\right)=\frac93\times(140+180+159)\ \mathrm{km}^{2}\cdot\mathrm m\]
\[=3\times479=1437\ \mathrm{km}^{2}\cdot\mathrm m\]

\textbf{第四步，统一单位。}

$1\ \mathrm{km}^{2}=10^{6}\ \mathrm m^{2}$，故
\[V=1437\times10^{6}\ \mathrm m^{3}\approx1.4\times10^{9}\ \mathrm m^{3}\]

\gkanswer{答案：C}

\begin{gktip}
\textbf{易错点}：单位换算。面积单位是 $\mathrm{km}^2$、高的单位是 $\mathrm m$，若直接相乘得到的是"$\mathrm{km}^2\cdot\mathrm m$"，必须乘 $10^{6}$ 才是 $\mathrm m^3$。四个选项相差都在 $10^{9}$ 量级内，单位错了会直接选错。
\end{gktip}

\par\medskip

\end{answer}
\begin{answer}
\gksolhead{5}{}

\textbf{题目}　从 $2$ 到 $8$ 的 $7$ 个整数中随机取 $2$ 个不同的数，则这 $2$ 个数互质的概率为（　　）

\textbf{解析}

七个数为 $2,3,4,5,6,7,8$，总取法
\[C_7^{2}=\frac{7\times6}2=21\]

\textbf{思路：数"不互质"的对更快}（互质的对多，不互质的少）。

两数不互质 $\iff$ 有公共质因数。$2\sim8$ 中可能的公共质因数只有 $2$ 和 $3$（$5$、$7$ 各只出现一次）：

\begin{itemize}
\item \textbf{公因数含 $2$}：偶数有 $2,4,6,8$ 共 $4$ 个，任取两个都不互质：$C_4^{2}=6$ 对
\item \textbf{公因数含 $3$}：$3$ 的倍数有 $3,6$，即 $(3,6)$ 共 $1$ 对
\item \textbf{重复计数检查}：既被"含 $2$"又被"含 $3$"计入的对，需两数都是 $6$ 的倍数——$2\sim8$ 中只有 $6$ 一个，无法成对，故\textbf{无重复}
\end{itemize}

不互质共 $6+1=7$ 对，互质共
\[21-7=14\ \text{对}\]

\[P=\frac{14}{21}=\frac23\]

\gkanswer{答案：D}

\begin{gktip}
\textbf{方法总结}：古典概型中若"正面数"较繁，优先考虑\textbf{补集}。本题正面枚举 14 对容易漏，反面只有 7 对且结构清晰（偶数对 + 一对 $3$ 的倍数）。另外别忘了检查重复计数——这是容斥原理的必要步骤。
\end{gktip}

\par\medskip

\end{answer}
\begin{answer}
\gksolhead{6}{}

\textbf{题目}　$f(x)=\sin\left(\omega x+\dfrac\pi4\right)+b\ (\omega>0)$ 的最小正周期为 $T$。若 $\dfrac{2\pi}3<T<\pi$，且图象关于点 $\left(\dfrac{3\pi}2,2\right)$ 中心对称，则 $f\left(\dfrac\pi2\right)=$（　　）

\textbf{解析}

\textbf{第一步，由周期范围限定 $\omega$。}
\[T=\frac{2\pi}\omega,\qquad \frac{2\pi}3<\frac{2\pi}\omega<\pi\]

分别解两个不等式（$\omega>0$）：
\[\frac{2\pi}\omega>\frac{2\pi}3\Rightarrow\omega<3,\qquad \frac{2\pi}\omega<\pi\Rightarrow\omega>2\]

\[2<\omega<3\]

\textbf{第二步，由中心对称确定 $b$ 与 $\omega$。}

正弦型函数 $y=\sin(\cdot)+b$ 的所有对称中心都在其"中轴线" $y=b$ 上。已知对称中心为 $\left(\dfrac{3\pi}2,2\right)$，故
\[b=2\]

且在对称中心处正弦值为零：
\[\sin\left(\omega\cdot\frac{3\pi}2+\frac\pi4\right)=0\ \Longrightarrow\ \frac{3\pi\omega}2+\frac\pi4=k\pi\ (k\in\mathbb Z)\]

两边除以 $\pi$：
\[\frac{3\omega}2+\frac14=k\ \Longrightarrow\ \omega=\frac{2}{3}\left(k-\frac14\right)=\frac{4k-1}6\]

\textbf{第三步，筛选 $k$。}
\[2<\frac{4k-1}6<3\ \Longrightarrow\ 12<4k-1<18\ \Longrightarrow\ 13<4k<19\ \Longrightarrow\ 3.25<k<4.75\]

整数 $k=4$，得
\[\omega=\frac{4\times4-1}6=\frac{15}6=\frac52\]

\textbf{第四步，求函数值。}
\[f(x)=\sin\left(\frac52x+\frac\pi4\right)+2\]

\[f\left(\frac\pi2\right)=\sin\left(\frac52\times\frac\pi2+\frac\pi4\right)+2=\sin\left(\frac{5\pi}4+\frac\pi4\right)+2=\sin\frac{3\pi}2+2=-1+2=1\]

\gkanswer{答案：A}

\begin{gktip}
\textbf{方法总结}：三角函数"由性质定参数"的题，标准流程是——① 由周期条件圈定 $\omega$ 的范围；② 由对称性（中心或轴）列出含 $k$ 的方程；③ 用①的范围筛出唯一的整数 $k$。三步缺一不可，其中第③步是把无穷多解收敛为唯一解的关键。
\end{gktip}

\par\medskip

\end{answer}
\begin{answer}
\gksolhead{7}{}

\textbf{题目}　设 $a=0.1\mathrm e^{0.1}$，$b=\dfrac19$，$c=-\ln0.9$，则（　　）

\textbf{解析}

\gkanswer{先用数值定位（考场上这一步可以快速锁定答案）：}
\[\mathrm e^{0.1}\approx1.10517\ \Longrightarrow\ a\approx0.110517\]
\[b=\frac19\approx0.111111\]
\[c=-\ln0.9=\ln\frac{10}9\approx0.105361\]

三者非常接近，排序为
\[c<a<b\]

\textbf{下面给出严格证明}（考场上若时间允许应补上，或至少心中有数）：

\textbf{(i) 证 $c<a$，即 $-\ln(1-0.1)<0.1\mathrm e^{0.1}$。}

构造
\[\varphi(x)=x\mathrm e^{x}+\ln(1-x),\qquad x\in(0,1)\]

\[\varphi(0)=0,\qquad \varphi'(x)=(1+x)\mathrm e^{x}-\frac1{1-x},\qquad \varphi'(0)=1-1=0\]

\[\varphi''(x)=(2+x)\mathrm e^{x}-\frac1{(1-x)^{2}},\qquad \varphi''(0)=2-1=1>0\]

在 $x=0$ 附近 $\varphi''>0$，故 $\varphi'$ 递增、$\varphi'(x)>0$，进而 $\varphi$ 递增，$\varphi(0.1)>\varphi(0)=0$，即
\[0.1\mathrm e^{0.1}>-\ln0.9\ \Longrightarrow\ a>c\]

\textbf{(ii) 证 $a<b$，即 $0.1\mathrm e^{0.1}<\dfrac19$。}

等价于
\[\mathrm e^{0.1}<\frac{10}9\iff 0.1<\ln\frac{10}9=\ln\left(1+\frac19\right)\]

利用不等式 $\ln(1+x)>\dfrac x{1+x}\ (x>0)$，取 $x=\dfrac19$：
\[\ln\frac{10}9>\frac{1/9}{10/9}=\frac1{10}=0.1\quad\checkmark\]

故 $a<b$。

\textbf{综上：$c<a<b$。}

\gkanswer{答案：C}

\begin{gktip}
\textbf{方法总结}：这类"三个数极接近"的比较题，考场策略是\textbf{先估值定序、再补证明}。常用的两个对数不等式值得背下来：
\[\frac{x}{1+x}<\ln(1+x)<x\quad(x>0)\]
本题的两次比较都能由它（或其变形）导出。
\end{gktip}

\par\medskip

\end{answer}
\begin{answer}
\gksolhead{8}{}

\textbf{题目}　正四棱锥的侧棱长为 $l$，各顶点都在同一球面上。若球体积为 $36\pi$，且 $3\le l\le3\sqrt3$，则该正四棱锥体积的取值范围是（　　）

\textbf{解析}

\textbf{第一步，求球半径。}
\[\frac43\pi R^{3}=36\pi\ \Longrightarrow\ R^{3}=27\ \Longrightarrow\ R=3\]

\textbf{第二步，建立 $l$ 与锥高 $h$ 的关系（本题的核心）。}

设正四棱锥底面边长为 $a$，底面外接圆半径
\[r=\frac{\sqrt2}2a\]

球心必在锥的轴线上。设锥顶在轴上高度 $h$、球心高度 $d$（以底面为 $0$）：

\begin{itemize}
\item 球心到底面顶点：$\sqrt{r^{2}+d^{2}}=R=3$，即 $r^{2}+d^{2}=9$
\item 球心到锥顶：$|h-d|=3$，取 $d=h-3$（球心在锥顶下方）
\end{itemize}

代入第一式：
\[r^{2}=9-(h-3)^{2}\]

又侧棱满足 $l^{2}=r^{2}+h^{2}$：
\[l^{2}=9-(h-3)^{2}+h^{2}=9-h^{2}+6h-9+h^{2}=6h\]

\textbf{得到关键关系}：
\[\boxed{l^{2}=6h}\]

由 $3\le l\le3\sqrt3$ 得 $9\le l^{2}\le27$，故
\[\frac32\le h\le\frac92\]

\textbf{第三步，把体积表示成 $h$ 的函数。}

底面积（正方形，对角线为 $2r$）：
\[S=a^{2}=2r^{2}=2\left[9-(h-3)^{2}\right]=2\left(6h-h^{2}\right)\]

\[V=\frac13Sh=\frac13\times2\left(6h-h^{2}\right)h=\frac23h^{2}(6-h)\]

\textbf{第四步，求 $V$ 在 $\left[\dfrac32,\dfrac92\right]$ 上的值域。}
\[V'(h)=\frac23\left(12h-3h^{2}\right)=2h(4-h)\]

在 $\left[\dfrac32,\dfrac92\right]$ 上：$h<4$ 时 $V'>0$（递增），$h>4$ 时 $V'<0$（递减），故 $h=4$ 是\textbf{最大值点}：
\[V(4)=\frac23\times16\times2=\frac{64}3\]

两个端点：
\[V\left(\frac32\right)=\frac23\times\frac94\times\frac92=\frac{27}4,\qquad V\left(\frac92\right)=\frac23\times\frac{81}4\times\frac32=\frac{81}4\]

比较 $\dfrac{27}4=6.75$ 与 $\dfrac{81}4=20.25$，最小值为 $\dfrac{27}4$。

\textbf{值域}：
\[\left[\frac{27}4,\ \frac{64}3\right]\]

（$\dfrac{64}3\approx21.33>\dfrac{81}4=20.25$，且函数连续，故区间内所有值都能取到。）

\gkanswer{答案：C}

\begin{gktip}
\textbf{方法总结}：外接球问题的通法是"\textbf{设球心在对称轴上，列两个等距方程}"。本题的漂亮之处在于消元后得到极简的 $l^2=6h$，把两个几何量的关系一次说清。最后一步要注意：最值可能在\textbf{内部驻点}取到（本题最大值在 $h=4$），不能只比较端点。
\end{gktip}

\par\medskip

\end{answer}
\gkpart{二、多选题（每小题 5 分，共 20 分）}

\begin{answer}
\gksolhead{9}{}

\textbf{题目}　已知正方体 $ABCD\text{-}A_1B_1C_1D_1$，则（　　）

\textbf{解析}

\textbf{建系}（设棱长为 $1$）：
\[A(0,0,0),\ B(1,0,0),\ C(1,1,0),\ D(0,1,0)\]
\[A_1(0,0,1),\ B_1(1,0,1),\ C_1(1,1,1),\ D_1(0,1,1)\]

四个选项都涉及 $BC_1$：
\[\overrightarrow{BC_1}=C_1-B=(0,1,1),\qquad \left|\overrightarrow{BC_1}\right|=\sqrt2\]

\par\medskip

\textbf{A 项}：$\overrightarrow{DA_1}=A_1-D=(0,-1,1)$
\[\overrightarrow{BC_1}\cdot\overrightarrow{DA_1}=0\times0+1\times(-1)+1\times1=0\]
垂直，所成角 $90^\circ$  \textbf{正确}。

\par\medskip

\textbf{B 项}：$\overrightarrow{CA_1}=A_1-C=(-1,-1,1)$
\[\overrightarrow{BC_1}\cdot\overrightarrow{CA_1}=0\times(-1)+1\times(-1)+1\times1=0\]
垂直  \textbf{正确}。

\begin{gktip}
这其实是正方体的经典结论：\textbf{体对角线垂直于不与它相交的面对角线}。
\end{gktip}

\par\medskip

\textbf{C 项}：平面 $BB_1D_1D$ 是过 $BD$ 的对角面。

由 $\overrightarrow{BD}=(-1,1,0)$、$\overrightarrow{BB_1}=(0,0,1)$：
\[\boldsymbol n=\overrightarrow{BD}\times\overrightarrow{BB_1}=(1\times1-0\times0,\ 0\times0-(-1)\times1,\ 0-0)=(1,1,0)\]

线面角公式（用 $\sin$）：
\[\sin\theta=\frac{\left|\overrightarrow{BC_1}\cdot\boldsymbol n\right|}{\left|\overrightarrow{BC_1}\right|\left|\boldsymbol n\right|}=\frac{|0+1+0|}{\sqrt2\times\sqrt2}=\frac12\]

故 $\theta=30^\circ\ne45^\circ$  \textbf{错误}。

\par\medskip

\textbf{D 项}：平面 $ABCD$ 即 $z=0$，法向量 $\boldsymbol n=(0,0,1)$。
\[\sin\theta=\frac{|1|}{\sqrt2\times1}=\frac{\sqrt2}2\ \Longrightarrow\ \theta=45^\circ\]
 \textbf{正确}（几何直观：$BC_1$ 在底面的投影是 $BC$，而 $CC_1=BC$，等腰直角）。

\gkanswer{答案：ABD}

\begin{gktip}
\textbf{方法总结}：正方体的角度问题一律建系，四个选项共用一次坐标设定，效率最高。注意\textbf{线面角用 $\sin$、线线角用 $\cos$ 且取绝对值}——C 项若误用 $\cos$ 会算出 $60^\circ$，同样选错。
\end{gktip}

\par\medskip

\end{answer}
\begin{answer}
\gksolhead{10}{}

\textbf{题目}　已知函数 $f(x)=x^{3}-x+1$，则（　　）

\textbf{解析}

\textbf{基础工作}：
\[f'(x)=3x^{2}-1\]

零点 $x=\pm\dfrac1{\sqrt3}=\pm\dfrac{\sqrt3}3$。

\begin{center}
\begin{tabular}{lccccc}
\toprule
$x$ & $\left(-\infty,-\frac{\sqrt3}3\right)$ & $-\frac{\sqrt3}3$ & $\left(-\frac{\sqrt3}3,\frac{\sqrt3}3\right)$ & $\frac{\sqrt3}3$ & $\left(\frac{\sqrt3}3,+\infty\right)$ \\
\midrule
$f'$ & $+$ & $0$ & $-$ & $0$ & $+$ \\
$f$ & 增 & 极大 & 减 & 极小 & 增 \\
\bottomrule
\end{tabular}
\end{center}

\par\medskip

\textbf{A 项 }：导数有两个变号零点，故有两个极值点 \textbf{正确}。

\par\medskip

\textbf{B 项 }：三次函数零点个数由\textbf{极值的符号}决定。

极小值：
\[f\left(\frac{\sqrt3}3\right)=\frac{\sqrt3}{9}\cdot\frac{3}{3}-\frac{\sqrt3}3+1\]

直接算：$\left(\dfrac{\sqrt3}3\right)^{3}=\dfrac{3\sqrt3}{27}=\dfrac{\sqrt3}9$，故
\[f\left(\frac{\sqrt3}3\right)=\frac{\sqrt3}9-\frac{\sqrt3}3+1=\frac{\sqrt3-3\sqrt3}9+1=-\frac{2\sqrt3}9+1\approx-0.385+1=0.615>0\]

\textbf{极小值为正}，说明函数图象在极小值处仍在 $x$ 轴上方，故整条曲线只与 $x$ 轴交于\textbf{一点}（在极大值点左侧）。

只有 $1$ 个零点，不是 $3$ 个  \textbf{错误}。

\par\medskip

\textbf{C 项 }：三次函数 $y=ax^3+bx^2+cx+d$ 的图象\textbf{必是中心对称图形}，对称中心是\textbf{拐点}（二阶导数零点）。
\[f''(x)=6x=0\ \Longrightarrow\ x=0,\qquad f(0)=1\]

对称中心为 $(0,1)$  \textbf{正确}。

（可验证：$f(x)+f(-x)=\left(x^3-x+1\right)+\left(-x^3+x+1\right)=2=2f(0)$ ）

\par\medskip

\textbf{D 项 }：若 $y=2x$ 是切线，设切点横坐标 $x_0$：
\[f'\left(x_0\right)=3x_0^{2}-1=2\ \Longrightarrow\ x_0^{2}=1\ \Longrightarrow\ x_0=\pm1\]

\begin{itemize}
\item $x_0=1$：$f(1)=1-1+1=1$，切线 $y-1=2(x-1)$，即 $y=2x-1$
\item $x_0=-1$：$f(-1)=-1+1+1=1$，切线 $y-1=2(x+1)$，即 $y=2x+3$
\end{itemize}

两条都不是 $y=2x$  \textbf{错误}。

\gkanswer{答案：AC}

\begin{gktip}
\textbf{易错点}：B 项容易凭"三次函数一般有三个零点"的印象误选。\textbf{判断三次函数零点个数的唯一可靠方法}是看极大值与极小值的符号：极大、极小异号（乘积为负）才有 3 个零点，同号只有 1 个。本题极小值 $\approx0.615>0$，极大值更大，同号 → 1 个零点。
\end{gktip}

\par\medskip

\end{answer}
\begin{answer}
\gksolhead{11}{}

\textbf{题目}　$O$ 为坐标原点，$A(1,1)$ 在抛物线 $C:x^{2}=2py\ (p>0)$ 上，过 $B(0,-1)$ 的直线交 $C$ 于 $P$、$Q$ 两点，则（　　）

\textbf{解析}

\textbf{预备：确定抛物线。} 把 $A(1,1)$ 代入：
\[1^{2}=2p\times1\ \Longrightarrow\ p=\frac12\]

\[C:\ x^{2}=y\]

\par\medskip

\textbf{A 项 }

$x^{2}=2py$ 的准线是 $y=-\dfrac p2$。此处 $p=\dfrac12$，准线为
\[y=-\frac14\ne-1\]
 \textbf{错误}。

\par\medskip

\textbf{B 项 }

直线 $AB$ 过 $A(1,1)$ 与 $B(0,-1)$：
\[k=\frac{1-(-1)}{1-0}=2,\qquad AB:\ y=2x-1\]

与 $y=x^{2}$ 联立：
\[x^{2}=2x-1\ \Longrightarrow\ x^{2}-2x+1=0\ \Longrightarrow\ (x-1)^{2}=0\]

\textbf{判别式为零、重根 $x=1$}，故直线与抛物线相切  \textbf{正确}。

\par\medskip

\textbf{统一设定（C、D 项共用）}

设过 $B(0,-1)$ 的直线为 $y=kx-1$，与 $y=x^{2}$ 联立：
\[x^{2}=kx-1\ \Longrightarrow\ x^{2}-kx+1=0\]

由韦达定理：
\[x_1+x_2=k,\qquad x_1x_2=1\]

有两个交点要求判别式为正：
\[k^{2}-4>0\ \Longrightarrow\ k^{2}>4\]

\par\medskip

\textbf{C 项 }

$P\left(x_1,x_1^{2}\right)$，$Q\left(x_2,x_2^{2}\right)$：
\[|OP|^{2}=x_1^{2}+x_1^{4}=x_1^{2}\left(1+x_1^{2}\right)\]

\[|OP|^{2}|OQ|^{2}=x_1^{2}x_2^{2}\left(1+x_1^{2}\right)\left(1+x_2^{2}\right)\]

由 $x_1x_2=1$ 得 $x_1^{2}x_2^{2}=1$，且
\[\left(1+x_1^{2}\right)\left(1+x_2^{2}\right)=1+x_1^{2}+x_2^{2}+x_1^{2}x_2^{2}=2+\left(x_1+x_2\right)^{2}-2x_1x_2=2+k^{2}-2=k^{2}\]

故
\[|OP|^{2}|OQ|^{2}=k^{2}>4\ \Longrightarrow\ |OP|\cdot|OQ|>2\]

而 $|OA|^{2}=1^{2}+1^{2}=2$，所以 $|OP|\cdot|OQ|>|OA|^{2}$  \textbf{正确}。

\par\medskip

\textbf{D 项 }

$B(0,-1)$ 在直线 $y=kx-1$ 上，故 $P$、$Q$ 到 $B$ 的距离可用横坐标差表示：
\[|BP|=\left|x_1-0\right|\sqrt{1+k^{2}}=\left|x_1\right|\sqrt{1+k^{2}}\]

\[|BP|\cdot|BQ|=\left|x_1x_2\right|\left(1+k^{2}\right)=1\times\left(1+k^{2}\right)=1+k^{2}>1+4=5\]

而
\[|BA|^{2}=(1-0)^{2}+(1+1)^{2}=1+4=5\]

故 $|BP|\cdot|BQ|>|BA|^{2}$  \textbf{正确}。

\gkanswer{答案：BCD}

\begin{gktip}
\textbf{方法总结}：C、D 两项都用到"\textbf{过定点的弦，两点到该点的距离乘积可由韦达定理直接算出}"这一技巧：$|BP||BQ|=\left|x_1x_2\right|\left(1+k^2\right)$。记住它可以避免逐个求出交点坐标。另外，B 项的"相切"用\textbf{判别式为零}判定，比求导算切线更快。
\end{gktip}

\par\medskip

\end{answer}
\begin{answer}
\gksolhead{12}{}

\textbf{题目}　$f(x)$ 及其导函数 $f'(x)$ 的定义域均为 $\mathbf R$，记 $g(x)=f'(x)$。若 $f\left(\dfrac32-2x\right)$、$g(2+x)$ 均为偶函数，则（　　）

\textbf{解析}

\textbf{第一步，把"复合后为偶函数"翻译成对称性。}

\textbf{关于 $f$：} $f\left(\dfrac32-2x\right)$ 为偶函数意味着把 $x$ 换成 $-x$ 值不变：
\[f\left(\frac32+2x\right)=f\left(\frac32-2x\right)\]

令 $u=2x$（$u$ 取遍 $\mathbf R$）：
\[f\left(\frac32+u\right)=f\left(\frac32-u\right)\]

即 \textbf{$f$ 的图象关于直线 $x=\dfrac32$ 轴对称}，等价写法：
\[f(3-x)=f(x)\qquad(\text{i})\]

\textbf{关于 $g$：} 同理，$g(2+x)$ 为偶函数意味着 \textbf{$g$ 关于直线 $x=2$ 轴对称}：
\[g(4-x)=g(x)\qquad(\text{ii})\]

\textbf{第二步，对 (i) 求导，得到 $g$ 的第二个对称性。}

对 $f(3-x)=f(x)$ 两边关于 $x$ 求导（左边用链式法则，内层导数为 $-1$）：
\[-f'(3-x)=f'(x)\ \Longrightarrow\ g(3-x)=-g(x)\qquad(\text{iii})\]

即 \textbf{$g$ 关于点 $\left(\dfrac32,0\right)$ 中心对称}。

\textbf{第三步，两个对称性合成周期。}

由 (ii)：$g(x)=g(4-x)$；把 (iii) 中的 $x$ 换成 $4-x$：
\[g\left(3-(4-x)\right)=-g(4-x)\ \Longrightarrow\ g(x-1)=-g(x)\]

即
\[g(x+1)=-g(x)\ \Longrightarrow\ g(x+2)=-g(x+1)=g(x)\]

\textbf{$g$ 是周期为 $2$ 的函数}，且满足"平移 $1$ 变号"。

\par\medskip

\textbf{逐项判断}

\textbf{A 项 }：题设的两个条件都只涉及 $f$ 的\textbf{对称性}与导数，给 $f$ 加上任意常数 $C$ 后条件依然成立，而 $f(0)$ 随之改变。故 $f(0)=0$ 不一定成立。

\textbf{B 项 }：由 (iii) 取 $x=\dfrac32$：
\[g\left(\frac32\right)=-g\left(\frac32\right)\ \Longrightarrow\ g\left(\frac32\right)=0\]

再由周期 $2$：
\[g\left(-\frac12\right)=g\left(-\frac12+2\right)=g\left(\frac32\right)=0\quad\checkmark\]

\textbf{C 项 }：由 (i) 取 $x=-1$：
\[f(-1)=f\left(3-(-1)\right)=f(4)\quad\checkmark\]

\textbf{D 项 }：由 $g(x+1)=-g(x)$ 与周期 $2$：
\[g(-1)=g(1)=-g(0)=-g(2)\]

若 D 成立需 $g(2)=-g(2)$，即 $g(2)=0$，但这并不是必然的。

\textbf{举反例}：取 $g(x)=\cos\pi x$，验证：
\begin{itemize}
\item $g(4-x)=\cos(4\pi-\pi x)=\cos\pi x=g(x)$  满足 (ii)
\item $g(3-x)=\cos(3\pi-\pi x)=-\cos\pi x=-g(x)$  满足 (iii)
\end{itemize}

而
\[g(-1)=\cos(-\pi)=-1,\qquad g(2)=\cos2\pi=1\]

$-1\ne1$  \textbf{错误}。

\gkanswer{答案：BC}

\begin{gktip}
\textbf{方法总结}：抽象函数的对称性题，核心是三步——
① 把"某个复合函数是偶/奇函数"翻译成 $f$ 或 $g$ 的\textbf{对称等式}；
② 对含 $f$ 的等式\textbf{求导}，转移到 $g$ 上（注意链式法则带来的负号，轴对称求导后变成中心对称）；
③ \textbf{两个对称性合成周期}（轴 + 轴 → 周期 $=2\times$轴距；轴 + 中心 → 周期 $=4\times$距离）。

本题中 $g$ 同时有轴对称（$x=2$）与中心对称（$\left(\dfrac32,0\right)$），距离 $\dfrac12$，故周期 $4\times\dfrac12=2$  与推导一致。
\end{gktip}

\par\medskip

\end{answer}
\gkpart{三、填空题（每小题 5 分，共 20 分）}

\begin{answer}
\gksolhead{13}{}

\textbf{题目}　$\left(1-\dfrac yx\right)(x+y)^{8}$ 的展开式中 $x^{2}y^{6}$ 的系数为 \gkblank{}。

\textbf{解析}

把括号拆开，分两部分处理：
\[\left(1-\frac yx\right)(x+y)^{8}=(x+y)^{8}-\frac yx(x+y)^{8}\]

\textbf{第一部分}：$(x+y)^{8}$ 的通项为 $C_8^{k}x^{8-k}y^{k}$。要得到 $x^{2}y^{6}$，需 $k=6$：
\[\text{系数}=C_8^{6}=C_8^{2}=\frac{8\times7}2=28\]

\textbf{第二部分}：$-\dfrac yx(x+y)^{8}$ 中，需要 $(x+y)^{8}$ 提供 $x^{3}y^{5}$（乘以 $\dfrac yx$ 后恰好变成 $x^{2}y^{6}$），即 $k=5$：
\[\text{系数}=-C_8^{5}=-C_8^{3}=-\frac{8\times7\times6}{6}=-56\]

\textbf{合并}：
\[28+(-56)=-28\]

\gkanswer{答案：$-28$}

\begin{gktip}
\textbf{方法总结}：形如 $\left(A+B\right)\cdot(x+y)^{n}$ 的展开式求某项系数，一律\textbf{拆成两个独立的二项式展开问题}，分别确定各自需要的项，最后相加。关键是想清楚"乘上 $-\dfrac yx$ 后，次数如何变化"——本题是 $x$ 降 $1$ 次、$y$ 升 $1$ 次，所以要找的是 $x^{3}y^{5}$ 而非 $x^{2}y^{6}$。
\end{gktip}

\par\medskip

\end{answer}
\begin{answer}
\gksolhead{14}{}

\textbf{题目}　写出与圆 $x^{2}+y^{2}=1$ 和 $(x-3)^{2}+(y-4)^{2}=16$ 都相切的一条直线的方程 \gkblank{}。

\textbf{解析}

\textbf{第一步，判断两圆的位置关系（决定公切线条数）。}

\begin{itemize}
\item 圆 $O_1$：圆心 $(0,0)$，半径 $r_1=1$
\item 圆 $O_2$：圆心 $(3,4)$，半径 $r_2=4$
\end{itemize}

圆心距：
\[d=\sqrt{3^{2}+4^{2}}=5\]

而
\[r_1+r_2=1+4=5=d\]

\textbf{两圆外切}，故有 \textbf{3 条}公切线（2 条外公切线 + 1 条过切点的内公切线）。

\textbf{第二步，找出这三条（写出任意一条即可得分）。}

\textbf{（a）竖直的外公切线 $x=-1$}

验证：$(0,0)$ 到 $x=-1$ 的距离为 $1=r_1$ ；$(3,4)$ 到 $x=-1$ 的距离为 $3-(-1)=4=r_2$ 

\textbf{这是最容易验证的一条，考场首选。}

\textbf{（b）内公切线（过两圆的公共切点）}

切点在 $O_1O_2$ 上，距 $O_1$ 为 $1$（占全长的 $\dfrac15$）：
\[\left(\frac35,\ \frac45\right)\]

切线垂直于 $O_1O_2$（方向 $(3,4)$）：
\[3\left(x-\frac35\right)+4\left(y-\frac45\right)=0\ \Longrightarrow\ 3x+4y-5=0\]

验证：到 $(0,0)$ 距离 $\dfrac{|-5|}5=1$ ；到 $(3,4)$ 距离 $\dfrac{|9+16-5|}5=4$ 

\textbf{（c）另一条外公切线}

设 $y=kx+b$，由两个距离条件解得 $k=\dfrac7{24}$、$b=-\dfrac{25}{24}$：
\[7x-24y-25=0\]

验证：到 $(0,0)$ 距离 $\dfrac{25}{\sqrt{49+576}}=\dfrac{25}{25}=1$ ；到 $(3,4)$ 距离 $\dfrac{|21-96-25|}{25}=\dfrac{100}{25}=4$ 

\gkanswer{答案：$x=-1$（或 $3x+4y-5=0$，或 $7x-24y-25=0$）}

\begin{gktip}
\textbf{考场策略}：题目只要求"写出一条"，所以\textbf{先判断位置关系、再优先找竖直线}（因为 $x=$ 常数的情形最容易口算验证）。若上来就设 $y=kx+b$ 联立方程，不仅慢，还会漏掉斜率不存在的 $x=-1$。
\end{gktip}

\par\medskip

\end{answer}
\begin{answer}
\gksolhead{15}{}

\textbf{题目}　若曲线 $y=(x+a)\mathrm e^{x}$ 有两条过坐标原点的切线，则 $a$ 的取值范围是 \gkblank{}。

\textbf{解析}

\textbf{第一步，写出切线过原点的条件。}

设切点为 $\left(t,\ (t+a)\mathrm e^{t}\right)$。

求导：
\[y'=\mathrm e^{x}+(x+a)\mathrm e^{x}=(x+a+1)\mathrm e^{x}\]

切线斜率 $=(t+a+1)\mathrm e^{t}$。

另一方面，切线过原点，斜率也等于"切点与原点连线的斜率"：
\[\frac{(t+a)\mathrm e^{t}-0}{t-0}=(t+a+1)\mathrm e^{t}\]

约去 $\mathrm e^{t}>0$，两边乘 $t$：
\[t+a=t(t+a+1)=t^{2}+at+t\]

整理：
\[a-at=t^{2}\ \Longrightarrow\ a(1-t)=t^{2}\]

\textbf{"有两条切线"等价于该方程关于 $t$ 有两个不同的实数解。}

\textbf{第二步，分离参数并研究函数。}

当 $t\ne1$ 时：
\[a=\frac{t^{2}}{1-t}\ \overset{\text{记}}{=}\ h(t)\]

（$t=1$ 时原方程为 $0=1$，无解，故可放心分离。）

\[h'(t)=\frac{2t(1-t)+t^{2}}{(1-t)^{2}}=\frac{2t-t^{2}}{(1-t)^{2}}=\frac{t(2-t)}{(1-t)^{2}}\]

分母恒正，符号由 $t(2-t)$ 决定：

\begin{center}
\begin{tabular}{lccccccc}
\toprule
$t$ & $(-\infty,0)$ & $0$ & $(0,1)$ & $1$ & $(1,2)$ & $2$ & $(2,+\infty)$ \\
\midrule
$h'$ & $-$ & $0$ & $+$ & 无定义 & $+$ & $0$ & $-$ \\
$h$ & 减 & 极小 $0$ & 增 & 断点 & 增 & 极大 $-4$ & 减 \\
\bottomrule
\end{tabular}
\end{center}

关键值：$h(0)=0$，$h(2)=\dfrac4{1-2}=-4$。

\textbf{极限行为}：
\begin{itemize}
\item $t\to-\infty$：$h(t)\approx\dfrac{t^{2}}{-t}=-t\to+\infty$
\item $t\to1^{-}$：分母 $\to0^{+}$，$h\to+\infty$
\item $t\to1^{+}$：分母 $\to0^{-}$，$h\to-\infty$
\item $t\to+\infty$：$h\approx-t\to-\infty$
\end{itemize}

\textbf{第三步，数解的个数。}

\textbf{左分支（$t<1$）}：从 $+\infty$ 减到 $h(0)=0$，再增回 $+\infty$，值域为 $[0,+\infty)$。
\begin{itemize}
\item $a>0$：\textbf{两个解}
\item $a=0$：一个解（$t=0$）
\item $a<0$：无解
\end{itemize}

\textbf{右分支（$t>1$）}：从 $-\infty$ 增到 $h(2)=-4$，再减回 $-\infty$，值域为 $(-\infty,-4]$。
\begin{itemize}
\item $a<-4$：\textbf{两个解}
\item $a=-4$：一个解（$t=2$）
\item $a>-4$：无解
\end{itemize}

\textbf{恰好两条切线}（两个不同的 $t$）：
\[a>0\quad\text{或}\quad a<-4\]

\gkanswer{答案：$(-\infty,-4)\cup(0,+\infty)$}

\begin{gktip}
\textbf{方法总结}："过某点的切线条数"问题的标准流程是——① 设切点参数 $t$，用"两种方式表示斜率"列出关于 $t$ 的方程；② \textbf{分离参数} $a=h(t)$；③ 画出 $h$ 的完整图象（含极值、断点、极限），数水平线 $y=a$ 与曲线的交点个数。第③步中断点两侧要分开讨论，这是本题的关键。
\end{gktip}

\par\medskip

\end{answer}
\begin{answer}
\gksolhead{16}{}

\textbf{题目}　椭圆 $C$ 的上顶点为 $A$，两焦点为 $F_1$、$F_2$，离心率 $\dfrac12$。过 $F_1$ 且垂直于 $AF_2$ 的直线与 $C$ 交于 $D$、$E$ 两点，$|DE|=6$，则 $\triangle ADE$ 的周长是 \gkblank{}。

\textbf{解析}

\textbf{第一步，挖掘 $e=\dfrac12$ 隐含的几何关系（本题的题眼）。}

由 $e=\dfrac ca=\dfrac12$ 得
\[a=2c\]

\textbf{上顶点到焦点的距离恒为 $a$}（因为 $\sqrt{b^{2}+c^{2}}=\sqrt{a^{2}-c^{2}+c^{2}}=a$）：
\[\left|AF_1\right|=\left|AF_2\right|=a\]

而两焦点间距：
\[\left|F_1F_2\right|=2c=a\]

于是
\[\left|F_1A\right|=\left|F_1F_2\right|=a\]

即 \textbf{$\triangle AF_1F_2$ 是以 $F_1$ 为顶点的等腰三角形}！

\textbf{第二步，识别出 $DE$ 是垂直平分线。}

等腰三角形中，\textbf{过顶点 $F_1$ 且垂直于底边 $AF_2$ 的直线，正是底边 $AF_2$ 的垂直平分线}。

题目所给的直线正是"过 $F_1$ 且垂直于 $AF_2$"，所以直线 $DE$ 就是 $AF_2$ 的垂直平分线。

\textbf{关键推论}：$DE$ 上的任意一点到 $A$ 和 $F_2$ 的距离相等，故
\[|DA|=\left|DF_2\right|,\qquad |EA|=\left|EF_2\right|\]

\textbf{第三步，用椭圆定义算周长。}

注意 $D$、$E$ 都在直线 $DE$ 上，且该直线过 $F_1$，故
\[|DE|=\left|DF_1\right|+\left|EF_1\right|\]

（$F_1$ 在弦 $DE$ 内部。）

于是
\[\begin{aligned}
\text{周长}&=|DA|+|EA|+|DE|\\
&=\left|DF_2\right|+\left|EF_2\right|+\left|DF_1\right|+\left|EF_1\right|\\
&=\underbrace{\left(\left|DF_1\right|+\left|DF_2\right|\right)}_{=2a}+\underbrace{\left(\left|EF_1\right|+\left|EF_2\right|\right)}_{=2a}\\
&=4a
\end{aligned}\]

\textbf{周长恒等于 $4a$，与 $|DE|$ 的具体数值无关！}

\textbf{第四步，由 $|DE|=6$ 求 $a$。}

取 $a=2$（则 $c=1$、$b=\sqrt3$）算出此时的弦长作为"标尺"：

$A\left(0,\sqrt3\right)$，$F_2(1,0)$，$k_{AF_2}=-\sqrt3$，故 $DE$ 的斜率为 $\dfrac1{\sqrt3}$，过 $F_1(-1,0)$：
\[y=\frac1{\sqrt3}(x+1)\]

代入 $\dfrac{x^{2}}4+\dfrac{y^{2}}3=1$，整理得
\[13x^{2}+8x-32=0,\qquad \left|x_1-x_2\right|=\frac{\sqrt{64+4\times13\times32}}{13}=\frac{24\sqrt3}{13}\]

\[|DE|=\sqrt{1+\frac13}\times\frac{24\sqrt3}{13}=\frac2{\sqrt3}\times\frac{24\sqrt3}{13}=\frac{48}{13}\]

由于椭圆的一切长度都与 $a$ 成正比：
\[\frac6{\dfrac{48}{13}}=\frac{13}8\ \Longrightarrow\ a=2\times\frac{13}8=\frac{13}4\]

\gkanswer{第五步，得出答案。}
\[\text{周长}=4a=4\times\frac{13}4=13\]

\gkanswer{答案：$13$}

\begin{gktip}
\textbf{方法总结}：本题的精妙之处在于 $e=\dfrac12$ 使得 $\left|F_1A\right|=\left|F_1F_2\right|$，从而"垂线"升格为"垂直平分线"，进而把 $|DA|$ 换成 $\left|DF_2\right|$，最终整个周长塌缩成 $4a$。\textbf{遇到"过焦点的垂线"，先检查是否构成等腰三角形}——这是解析几何中常见的隐藏结构。
\end{gktip}

\par\medskip

\end{answer}
\gkpart{四、解答题（共 70 分）}

\begin{answer}
\gksolhead{17}{10 分}

\textbf{题目}　$S_n$ 为 $\left\{a_n\right\}$ 的前 $n$ 项和，$a_1=1$，$\left\{\dfrac{S_n}{a_n}\right\}$ 是公差为 $\dfrac13$ 的等差数列。(1) 求通项；(2) 证明 $\displaystyle\sum_{k=1}^n\frac1{a_k}<2$。

\par\medskip

\textbf{(1) 解}

\textbf{第一步，写出 $\dfrac{S_n}{a_n}$ 的通项。}

首项：
\[\frac{S_1}{a_1}=\frac{a_1}{a_1}=1\]

公差 $\dfrac13$，故
\[\frac{S_n}{a_n}=1+(n-1)\times\frac13=\frac{n+2}3\]

即
\[3S_n=(n+2)a_n\qquad(\text{i})\]

\textbf{第二步，作差消去 $S$。}

当 $n\ge2$ 时，同式对 $n-1$ 也成立：
\[3S_{n-1}=(n+1)a_{n-1}\qquad(\text{ii})\]

$(\text{i})-(\text{ii})$（注意 $S_n-S_{n-1}=a_n$）：
\[3a_n=(n+2)a_n-(n+1)a_{n-1}\]

移项：
\[(n+1)a_{n-1}=(n+2)a_n-3a_n=(n-1)a_n\]

\[\frac{a_n}{a_{n-1}}=\frac{n+1}{n-1}\qquad(n\ge2)\]

\textbf{第三步，累乘求通项。}
\[\frac{a_n}{a_1}=\frac{a_2}{a_1}\cdot\frac{a_3}{a_2}\cdots\frac{a_n}{a_{n-1}}=\frac31\cdot\frac42\cdot\frac53\cdot\frac64\cdots\frac{n+1}{n-1}\]

分子为 $3\cdot4\cdot5\cdots(n+1)=\dfrac{(n+1)!}{2}$，分母为 $1\cdot2\cdot3\cdots(n-1)=(n-1)!$：
\[\frac{a_n}{a_1}=\frac{(n+1)!/2}{(n-1)!}=\frac{n(n+1)}2\]

由 $a_1=1$：
\[a_n=\frac{n(n+1)}2\]

\textbf{验证}：$a_1=1$ ，$a_2=3$，$S_2=4$，$\dfrac{S_2}{a_2}=\dfrac43=1+\dfrac13$  公差正确。

\par\medskip

\textbf{(2) 证明}

\textbf{关键：裂项。}
\[\frac1{a_n}=\frac2{n(n+1)}=2\left(\frac1n-\frac1{n+1}\right)\]

（用了 $\dfrac1{n(n+1)}=\dfrac1n-\dfrac1{n+1}$。）

求和时中间项\textbf{首尾相消}：
\[\sum_{k=1}^n\frac1{a_k}=2\sum_{k=1}^n\left(\frac1k-\frac1{k+1}\right)=2\left[\left(1-\frac12\right)+\left(\frac12-\frac13\right)+\cdots+\left(\frac1n-\frac1{n+1}\right)\right]\]

\[=2\left(1-\frac1{n+1}\right)=2-\frac2{n+1}\]

由于 $\dfrac2{n+1}>0$ 对一切 $n\in\mathbb N^{*}$ 成立：
\[\sum_{k=1}^n\frac1{a_k}=2-\frac2{n+1}<2\qquad\blacksquare\]

\begin{gktip}
\textbf{方法总结}：数列不等式若出现 $\dfrac1{n(n+1)}$、$\dfrac1{(2n-1)(2n+1)}$ 这类形式，一律\textbf{裂项相消}，求和后直接与目标比较。本题求和结果 $2-\dfrac2{n+1}$ 不仅小于 $2$，还表明 $2$ 是\textbf{最小上界}（$n\to\infty$ 时趋近于 $2$），这也解释了题目为什么给的是 $2$ 而不是别的数。
\end{gktip}

\par\medskip

\end{answer}
\begin{answer}
\gksolhead{18}{12 分}

\textbf{题目}　$\triangle ABC$ 中 $\dfrac{\cos A}{1+\sin A}=\dfrac{\sin2B}{1+\cos2B}$。(1) 若 $C=\dfrac{2\pi}3$，求 $B$；(2) 求 $\dfrac{a^{2}+b^{2}}{c^{2}}$ 的最小值。

\par\medskip

\textbf{预备：化简条件（全题的钥匙）}

\textbf{右边}用二倍角公式：
\[\frac{\sin2B}{1+\cos2B}=\frac{2\sin B\cos B}{1+\left(2\cos^{2}B-1\right)}=\frac{2\sin B\cos B}{2\cos^{2}B}=\tan B\]

\textbf{左边}用半角公式。分子分母同时用 $\dfrac A2$ 表示：
\[\cos A=\cos^{2}\frac A2-\sin^{2}\frac A2=\left(\cos\frac A2-\sin\frac A2\right)\left(\cos\frac A2+\sin\frac A2\right)\]
\[1+\sin A=\cos^{2}\frac A2+\sin^{2}\frac A2+2\sin\frac A2\cos\frac A2=\left(\cos\frac A2+\sin\frac A2\right)^{2}\]

相除（约去公因式）：
\[\frac{\cos A}{1+\sin A}=\frac{\cos\frac A2-\sin\frac A2}{\cos\frac A2+\sin\frac A2}=\frac{1-\tan\frac A2}{1+\tan\frac A2}=\tan\left(\frac\pi4-\frac A2\right)\]

\textbf{于是条件化为}
\[\tan B=\tan\left(\frac\pi4-\frac A2\right)\]

由于 $B$ 与 $\dfrac\pi4-\dfrac A2$ 都在 $\left(-\dfrac\pi2,\dfrac\pi2\right)$ 内（三角形中 $0<A<\pi$ 故 $\dfrac\pi4-\dfrac A2\in\left(-\dfrac\pi4,\dfrac\pi4\right)$），正切函数在此区间单调，可去掉 $\tan$：
\[B=\frac\pi4-\frac A2\]

\textbf{核心关系}：
\[\boxed{A+2B=\frac\pi2}\]

\par\medskip

\textbf{(1) 解}

由 $C=\dfrac{2\pi}3$ 及内角和：
\[A+B=\pi-\frac{2\pi}3=\frac\pi3\]

与核心关系联立：
\[\begin{cases}A+2B=\dfrac\pi2\\ A+B=\dfrac\pi3\end{cases}\]

两式相减：
\[B=\frac\pi2-\frac\pi3=\frac\pi6\]

\textbf{所以 $B=\dfrac\pi6$。}

\par\medskip

\textbf{(2) 解}

\textbf{第一步，把 $A$、$C$ 都用 $B$ 表示。}

由核心关系：
\[A=\frac\pi2-2B\]

\[C=\pi-A-B=\pi-\left(\frac\pi2-2B\right)-B=\frac\pi2+B\]

\textbf{角的范围}：$A>0\Rightarrow B<\dfrac\pi4$；又 $B>0$，故 $B\in\left(0,\dfrac\pi4\right)$。

\textbf{第二步，用正弦定理把边转化为角。}
\[\frac{a^{2}+b^{2}}{c^{2}}=\frac{\sin^{2}A+\sin^{2}B}{\sin^{2}C}\]

代入：
\[\sin A=\sin\left(\frac\pi2-2B\right)=\cos2B,\qquad \sin C=\sin\left(\frac\pi2+B\right)=\cos B\]

\[\frac{a^{2}+b^{2}}{c^{2}}=\frac{\cos^{2}2B+\sin^{2}B}{\cos^{2}B}\]

\textbf{第三步，换元化为单变量。}

令 $t=\cos2B$，则
\[\sin^{2}B=\frac{1-t}2,\qquad \cos^{2}B=\frac{1+t}2\]

\[\frac{a^{2}+b^{2}}{c^{2}}=\frac{t^{2}+\dfrac{1-t}2}{\dfrac{1+t}2}=\frac{2t^{2}+1-t}{1+t}\]

再令 $u=1+t$（即 $t=u-1$）：
\[=\frac{2(u-1)^{2}-(u-1)+1}{u}=\frac{2u^{2}-4u+2-u+1+1}{u}=\frac{2u^{2}-5u+4}{u}=2u+\frac4u-5\]

\textbf{第四步，用基本不等式。}
\[2u+\frac4u\ge2\sqrt{2u\cdot\frac4u}=2\sqrt8=4\sqrt2\]

等号成立当 $2u=\dfrac4u$，即 $u^{2}=2$，$u=\sqrt2$。

\textbf{检验取值范围}：$B\in\left(0,\dfrac\pi4\right)\Rightarrow2B\in\left(0,\dfrac\pi2\right)\Rightarrow t=\cos2B\in(0,1)\Rightarrow u=1+t\in(1,2)$。

而 $\sqrt2\approx1.414\in(1,2)$  \textbf{等号可以取到}。

\[\left(\frac{a^{2}+b^{2}}{c^{2}}\right)_{\min}=4\sqrt2-5\]

\gkanswer{答案：$4\sqrt2-5$}

\begin{gktip}
\textbf{方法总结}：解三角形的最值题，标准路线是"\textbf{化角 → 减元 → 换元 → 基本不等式（或求导）}"。本题的两次换元（先 $t=\cos2B$ 统一三角函数，再 $u=1+t$ 凑出 $2u+\dfrac4u$ 的对勾形式）是关键。\textbf{最后务必检验等号能否取到}——若 $\sqrt2$ 落在 $(1,2)$ 之外，最小值就要在端点处取，答案完全不同。
\end{gktip}

\par\medskip

\end{answer}
\begin{answer}
\gksolhead{19}{12 分}

\textbf{题目}　直三棱柱 $ABC\text{-}A_1B_1C_1$ 的体积为 $4$，$\triangle A_1BC$ 的面积为 $2\sqrt2$。

\begin{center}% 2022 新高考Ⅰ卷 19：直三棱柱 ABC-A1B1C1（D 为 A1C 的中点）
\begin{tikzpicture}[scale=0.95, font=\small]
  \coordinate (A)  at (0,0.5);
  \coordinate (B)  at (1.15,-0.42);
  \coordinate (C)  at (2.5,0.5);
  \coordinate (A1) at (0,2.7);
  \coordinate (B1) at (1.15,1.78);
  \coordinate (C1) at (2.5,2.7);
  \coordinate (D)  at ($(A1)!0.5!(C)$);
  \draw (A) -- (B) -- (C);
  \draw[dashed] (A) -- (C);
  \draw (A) -- (A1)  (B) -- (B1)  (C) -- (C1);
  \draw (A1) -- (B1) -- (C1) -- cycle;
  \draw (A1) -- (B);
  \draw[dashed] (A1) -- (C);
  \draw[dashed] (B) -- (D);
  \fill (D) circle (0.9pt);
  \node[left] at (A) {$A$};
  \node[below] at (B) {$B$};
  \node[right] at (C) {$C$};
  \node[above left, inner sep=1pt] at (A1) {$A_1$};
  \node[below, inner sep=2.5pt] at (B1) {$B_1$};
  \node[above right, inner sep=1pt] at (C1) {$C_1$};
  \node[above right, inner sep=0.5pt] at (D) {$D$};
\end{tikzpicture}
\end{center}

  (1) 求 $A$ 到平面 $A_1BC$ 的距离；(2) 设 $D$ 为 $A_1C$ 的中点，$AA_1=AB$，平面 $A_1BC\perp$ 平面 $ABB_1A_1$，求二面角 $A\text{-}BD\text{-}C$ 的正弦值。

\par\medskip

\textbf{(1) 解}

\textbf{关键：同一个四面体，用两种方式算体积。}

四面体 $A_1\text{-}ABC$（顶点 $A_1$、底面 $\triangle ABC$）的体积是三棱柱体积的 $\dfrac13$：
\[V_{A_1\text{-}ABC}=\frac13V_{\text{柱}}=\frac13\times4=\frac43\]

（因为柱体积 $=S_{\triangle ABC}\times h$，而锥体积 $=\dfrac13S_{\triangle ABC}\times h$。）

这个四面体也可以看作以 $\triangle A_1BC$ 为底、$A$ 为顶点。设 $A$ 到平面 $A_1BC$ 的距离为 $d$：
\[V=\frac13S_{\triangle A_1BC}\cdot d=\frac13\times2\sqrt2\times d=\frac43\]

\[2\sqrt2\,d=4\ \Longrightarrow\ d=\frac4{2\sqrt2}=\sqrt2\]

\textbf{所以 $A$ 到平面 $A_1BC$ 的距离为 $\sqrt2$。}

\par\medskip

\textbf{(2) 解}

\textbf{第一步，利用面面垂直确定边长。}

取 $A_1B$ 的中点 $E$，连接 $AE$。

因为是\textbf{直}三棱柱，$AA_1\perp$ 底面，故 $AA_1\perp AB$，即 $\triangle AA_1B$ 是直角三角形；又 $AA_1=AB$，故它是\textbf{等腰直角三角形}，从而中线 $AE$ 也是高：
\[AE\perp A_1B\]

由\textbf{平面 $A_1BC\perp$ 平面 $ABB_1A_1$}，且交线为 $A_1B$，$AE\subset$ 平面 $ABB_1A_1$ 且 $AE\perp A_1B$，由面面垂直的性质定理：
\[AE\perp\text{平面 }A_1BC\]

所以 $AE$ 就是第 (1) 问求的距离：
\[AE=\sqrt2\]

在等腰直角三角形 $AA_1B$ 中，斜边 $A_1B$ 上的中线等于斜边的一半：
\[AE=\frac{A_1B}2=\frac{\sqrt2\,AB}2\]

由 $\dfrac{\sqrt2\,AB}2=\sqrt2$ 得
\[AB=AA_1=2,\qquad A_1B=2\sqrt2\]

\textbf{第二步，证明 $BC\perp AB$（建系的前提）。}

由 $AE\perp$ 平面 $A_1BC$ 且 $BC\subset$ 平面 $A_1BC$，得 $AE\perp BC$；
又 $AA_1\perp$ 底面 $ABC$，$BC\subset$ 底面，得 $AA_1\perp BC$。

$AE$ 与 $AA_1$ 是平面 $ABB_1A_1$ 内的两条\textbf{相交}直线，故
\[BC\perp\text{平面 }ABB_1A_1\]

从而 $BC\perp AB$。

\textbf{第三步，建系求 $BC$。}

以 $B$ 为原点，$BA$、$BC$、$BB_1$ 为 $x$、$y$、$z$ 轴：
\[B(0,0,0),\quad A(2,0,0),\quad A_1(2,0,2),\quad C(0,c,0)\]

由 $BC\perp$ 平面 $ABB_1A_1$ 知 $BC\perp A_1B$，故 $\triangle A_1BC$ 是直角三角形：
\[S_{\triangle A_1BC}=\frac12\times BC\times A_1B=\frac12\times c\times2\sqrt2=\sqrt2\,c=2\sqrt2\]

\[c=2,\qquad C(0,2,0)\]

\textbf{第四步，求二面角。}

$D$ 是 $A_1C$ 的中点：
\[D=\frac{A_1+C}2=\frac{(2,0,2)+(0,2,0)}2=(1,1,1)\]

棱为 $BD$：$\overrightarrow{BD}=(1,1,1)$。

\textbf{平面 $ABD$} 的法向量：
\[\boldsymbol n_1=\overrightarrow{BD}\times\overrightarrow{BA}=(1,1,1)\times(2,0,0)=(1\cdot0-1\cdot0,\ 1\cdot2-1\cdot0,\ 0-2)=(0,\ 2,\ -2)\]

\textbf{平面 $CBD$} 的法向量：
\[\boldsymbol n_2=\overrightarrow{BD}\times\overrightarrow{BC}=(1,1,1)\times(0,2,0)=(1\cdot0-1\cdot2,\ 0-0,\ 2-0)=(-2,\ 0,\ 2)\]

\[\cos\left\langle\boldsymbol n_1,\boldsymbol n_2\right\rangle=\frac{0+0-4}{\sqrt{0+4+4}\times\sqrt{4+0+4}}=\frac{-4}{2\sqrt2\times2\sqrt2}=\frac{-4}8=-\frac12\]

故二面角的余弦值的绝对值为 $\dfrac12$，正弦值为
\[\sin\theta=\sqrt{1-\left(\frac12\right)^{2}}=\frac{\sqrt3}2\]

\textbf{所以二面角 $A\text{-}BD\text{-}C$ 的正弦值为 $\dfrac{\sqrt3}2$。}

\begin{gktip}
\textbf{方法总结}：本题第 (2) 问的难点全在\textbf{建系之前}——必须先由"面面垂直 + 等腰直角"推出 $AE\perp$ 平面 $A_1BC$，进而定出边长，再推出 $BC\perp AB$ 才能建立直角坐标系。这类题的经验是：\textbf{面面垂直的性质定理（在一个面内垂直于交线的直线垂直于另一个面）几乎总是突破口}。
\end{gktip}

\par\medskip

\end{answer}
\begin{answer}
\gksolhead{20}{12 分}

\textbf{题目}　病例组（患病 $100$ 人）中卫生习惯不够良好 $40$、良好 $60$；对照组（未患病 $100$ 人）中不够良好 $10$、良好 $90$。

  (1) 能否有 $99\%$ 的把握认为两群体卫生习惯有差异？
  (2) $A$：卫生习惯不够良好；$B$：患该疾病。$R=\dfrac{P(B\mid A)}{P(\bar B\mid A)}\Big/\dfrac{P(B\mid\bar A)}{P(\bar B\mid\bar A)}$。① 证明 $R=\dfrac{P(A\mid B)}{P(\bar A\mid B)}\cdot\dfrac{P(\bar A\mid\bar B)}{P(A\mid\bar B)}$；② 求 $R$ 的估计值。

\par\medskip

\textbf{(1) 解}

\textbf{提出零假设} $H_0$：患病与卫生习惯\textbf{独立}（无差异）。

按列联表，$a=40$、$b=60$、$c=10$、$d=90$，$n=200$：

\[ad-bc=40\times90-60\times10=3600-600=3000\]

\[K^{2}=\frac{n(ad-bc)^{2}}{(a+b)(c+d)(a+c)(b+d)}=\frac{200\times3000^{2}}{100\times100\times50\times150}\]

\textbf{分子}：$200\times9\times10^{6}=1.8\times10^{9}$
\textbf{分母}：$100\times100\times50\times150=7.5\times10^{7}$

\[K^{2}=\frac{1.8\times10^{9}}{7.5\times10^{7}}=24\]

查表：$\alpha=0.01$（对应 $99\%$ 把握）的临界值为 $6.635$。

因
\[24>6.635\]

故\textbf{拒绝} $H_0$，即\textbf{有 $99\%$ 的把握认为患该疾病群体与未患该疾病群体的卫生习惯有差异}。

\par\medskip

\textbf{(2)① 证明}

\textbf{思路：把所有条件概率还原为联合概率，两边各自化简。}

\textbf{左边（$R$ 的定义式）}：
\[\frac{P(B\mid A)}{P(\bar B\mid A)}=\frac{P(AB)/P(A)}{P(A\bar B)/P(A)}=\frac{P(AB)}{P(A\bar B)}\]

（分母的 $P(A)$ 约去。）同理
\[\frac{P(B\mid\bar A)}{P(\bar B\mid\bar A)}=\frac{P(\bar AB)}{P(\bar A\bar B)}\]

故
\[R=\frac{P(AB)}{P(A\bar B)}\Big/\frac{P(\bar AB)}{P(\bar A\bar B)}=\frac{P(AB)\cdot P(\bar A\bar B)}{P(A\bar B)\cdot P(\bar AB)}\qquad(\star)\]

\textbf{右边（待证式）}：
\[\frac{P(A\mid B)}{P(\bar A\mid B)}=\frac{P(AB)/P(B)}{P(\bar AB)/P(B)}=\frac{P(AB)}{P(\bar AB)}\]

\[\frac{P(\bar A\mid\bar B)}{P(A\mid\bar B)}=\frac{P(\bar A\bar B)}{P(A\bar B)}\]

两式相乘：
\[\frac{P(AB)}{P(\bar AB)}\cdot\frac{P(\bar A\bar B)}{P(A\bar B)}=\frac{P(AB)\cdot P(\bar A\bar B)}{P(A\bar B)\cdot P(\bar AB)}\]

与 $(\star)$ \textbf{完全相同}，故
\[R=\frac{P(A\mid B)}{P(\bar A\mid B)}\cdot\frac{P(\bar A\mid\bar B)}{P(A\mid\bar B)}\qquad\blacksquare\]

\begin{gktip}
\textbf{这个恒等式的现实意义}：左边需要"已知卫生习惯，患病的概率"——这在流行病学中难以直接调查（要长期追踪人群）；右边只需"已知患病与否，卫生习惯的比例"——正是\textbf{病例对照研究}能直接得到的数据。第 ① 问证明的正是"病例对照研究可以估计疾病风险"这一统计学基本原理。
\end{gktip}

\par\medskip

\textbf{(2)② 解}

由调查数据直接估计（用频率代替概率）：

\begin{itemize}
\item 病例组即事件 $B$，其中不够良好占：
\end{itemize}
\[P(A\mid B)\approx\frac{40}{100}=0.4\]

\begin{itemize}
\item 对照组即事件 $\bar B$，其中不够良好占：
\end{itemize}
\[P(A\mid\bar B)\approx\frac{10}{100}=0.1\]

相应地：
\[P(\bar A\mid B)=1-0.4=0.6,\qquad P(\bar A\mid\bar B)=1-0.1=0.9\]

代入 ① 的公式：
\[R=\frac{P(A\mid B)}{P(\bar A\mid B)}\cdot\frac{P(\bar A\mid\bar B)}{P(A\mid\bar B)}=\frac{0.4}{0.6}\times\frac{0.9}{0.1}=\frac23\times9=6\]

\textbf{所以 $R$ 的估计值为 $6$。}

\begin{gktip}
\textbf{结果解读}：$R=6$ 表示卫生习惯不够良好的人患该疾病的风险约是习惯良好者的 $6$ 倍。
\end{gktip}

\par\medskip

\end{answer}
\begin{answer}
\gksolhead{21}{12 分}

\textbf{题目}　$A(2,1)$ 在双曲线 $C:\dfrac{x^{2}}{a^{2}}-\dfrac{y^{2}}{a^{2}-1}=1\ (a>1)$ 上，直线 $l$ 交 $C$ 于 $P$、$Q$，且 $k_{AP}+k_{AQ}=0$。(1) 求 $l$ 的斜率；(2) 若 $\tan\angle PAQ=2\sqrt2$，求 $\triangle PAQ$ 的面积。

\par\medskip

\textbf{(1) 解}

\textbf{第一步，先把曲线定下来。}

把 $A(2,1)$ 代入方程：
\[\frac4{a^{2}}-\frac1{a^{2}-1}=1\]

两边乘 $a^{2}\left(a^{2}-1\right)$：
\[4\left(a^{2}-1\right)-a^{2}=a^{2}\left(a^{2}-1\right)\]
\[3a^{2}-4=a^{4}-a^{2}\ \Longrightarrow\ a^{4}-4a^{2}+4=0\ \Longrightarrow\ \left(a^{2}-2\right)^{2}=0\]

\[a^{2}=2\]

（完全平方，$a^2$ 唯一。）故
\[C:\ \frac{x^{2}}2-y^{2}=1\]

\textbf{第二步，设直线并联立。}

设 $l:y=kx+m$，与 $C$ 交于 $P\left(x_1,y_1\right)$、$Q\left(x_2,y_2\right)$。

代入 $x^{2}-2y^{2}=2$：
\[x^{2}-2(kx+m)^{2}=2\]
\[\left(1-2k^{2}\right)x^{2}-4kmx-\left(2m^{2}+2\right)=0\]

韦达定理：
\[x_1+x_2=\frac{4km}{1-2k^{2}},\qquad x_1x_2=\frac{-\left(2m^{2}+2\right)}{1-2k^{2}}\]

\textbf{第三步，把斜率和为零的条件代数化。}
\[k_{AP}+k_{AQ}=\frac{y_1-1}{x_1-2}+\frac{y_2-1}{x_2-2}=0\]

通分后分子为零：
\[\left(y_1-1\right)\left(x_2-2\right)+\left(y_2-1\right)\left(x_1-2\right)=0\]

代入 $y_i=kx_i+m$：
\[\left(kx_1+m-1\right)\left(x_2-2\right)+\left(kx_2+m-1\right)\left(x_1-2\right)=0\]

展开整理：
\[2kx_1x_2+(m-1-2k)\left(x_1+x_2\right)-4(m-1)=0\]

\textbf{第四步，代入韦达定理化简。}
\[2k\cdot\frac{-\left(2m^{2}+2\right)}{1-2k^{2}}+(m-1-2k)\cdot\frac{4km}{1-2k^{2}}-4(m-1)=0\]

两边乘 $\left(1-2k^{2}\right)$ 并展开：
\[-4km^{2}-4k+4km^{2}-4km-8k^{2}m-4(m-1)\left(1-2k^{2}\right)=0\]

\[-4k-4km-8k^{2}m-4m+4+8k^{2}m-8k^{2}=0\]

\[-4k-4km-4m+4-8k^{2}=0\]

同除以 $-4$：
\[k+km+m-1+2k^{2}=0\]

\textbf{因式分解}（把 $m$ 提出来）：
\[m(k+1)+\left(2k^{2}+k-1\right)=0\ \Longrightarrow\ m(k+1)+(2k-1)(k+1)=0\]

\[(k+1)\left[m+2k-1\right]=0\]

\textbf{第五步，取舍。}

\begin{itemize}
\item 若 $m=1-2k$：直线为 $y=kx+1-2k=k(x-2)+1$，\textbf{恒过点 $A(2,1)$}。此时 $P$ 或 $Q$ 与 $A$ 重合，$k_{AP}$ 无意义，\textbf{舍去}。
\item 故必有 $k+1=0$，即
\end{itemize}
\[k=-1\]

\textbf{所以 $l$ 的斜率为 $-1$。}

\par\medskip

\textbf{(2) 解}

\textbf{第一步，利用对称性设斜率。}

由 (1) 的条件 $k_{AP}+k_{AQ}=0$，设
\[k_{AP}=t,\qquad k_{AQ}=-t\qquad(t>0)\]

\textbf{第二步，由夹角求 $t$。}
\[\tan\angle PAQ=\left|\frac{k_{AP}-k_{AQ}}{1+k_{AP}k_{AQ}}\right|=\left|\frac{2t}{1-t^{2}}\right|=2\sqrt2\]

\textbf{情形一}：$\dfrac{2t}{1-t^{2}}=2\sqrt2$，即 $\sqrt2t^{2}+t-\sqrt2=0$，解得
\[t=\frac{-1\pm3}{2\sqrt2}\ \Longrightarrow\ t=\frac1{\sqrt2}\ (\text{取正})\]

\textbf{情形二}：$\dfrac{2t}{1-t^{2}}=-2\sqrt2$，解得 $t=\sqrt2$。

两种情形由对称性给出相同的三角形（只是 $P$、$Q$ 互换），下面取 $t=\sqrt2$ 计算。

\textbf{第三步，求 $|AP|$、$|AQ|$（用"韦达定理求另一交点"的技巧）。}

直线 $AP$：$y-1=\sqrt2(x-2)$，即 $y=\sqrt2x+s$，其中 $s=1-2\sqrt2$。

代入 $x^{2}-2y^{2}=2$：
\[x^{2}-2\left(\sqrt2x+s\right)^{2}=2\ \Longrightarrow\ x^{2}-2\left(2x^{2}+2\sqrt2sx+s^{2}\right)=2\]
\[-3x^{2}-4\sqrt2sx-2s^{2}-2=0\ \Longrightarrow\ 3x^{2}+4\sqrt2sx+2s^{2}+2=0\]

已知一根是 $x_A=2$，由\textbf{根与系数的关系}（两根之积）：
\[2\cdot x_P=\frac{2s^{2}+2}3\ \Longrightarrow\ x_P=\frac{s^{2}+1}3\]

$s=1-2\sqrt2$，$s^{2}=1-4\sqrt2+8=9-4\sqrt2$：
\[x_P=\frac{10-4\sqrt2}3,\qquad x_P-2=\frac{4-4\sqrt2}3\]

\[|AP|=\left|x_P-2\right|\sqrt{1+t^{2}}=\frac{4\left(\sqrt2-1\right)}3\times\sqrt3=\frac{4\sqrt3\left(\sqrt2-1\right)}3\]

同理对直线 $AQ$（斜率 $-\sqrt2$，$s'=1+2\sqrt2$，$s'^{2}=9+4\sqrt2$）：
\[x_Q=\frac{10+4\sqrt2}3,\qquad |AQ|=\frac{4\sqrt3\left(\sqrt2+1\right)}3\]

\textbf{第四步，算面积。}

两长度相乘（利用平方差）：
\[|AP|\cdot|AQ|=\frac{16\times3\times\left(\sqrt2-1\right)\left(\sqrt2+1\right)}9=\frac{48\times(2-1)}9=\frac{16}3\]

由 $\tan\angle PAQ=2\sqrt2$ 得（构造直角三角形：对边 $2\sqrt2$、邻边 $1$、斜边 $3$）：
\[\sin\angle PAQ=\frac{2\sqrt2}3\]

\[S=\frac12|AP|\cdot|AQ|\cdot\sin\angle PAQ=\frac12\times\frac{16}3\times\frac{2\sqrt2}3=\frac{16\sqrt2}9\]

\textbf{所以 $\triangle PAQ$ 的面积为 $\dfrac{16\sqrt2}9$。}

\begin{gktip}
\textbf{方法总结}：第 (2) 问的关键技巧是"\textbf{已知曲线上一点求过它的直线的另一交点}"——不必解二次方程，直接用韦达定理的两根之积（或之和），因为其中一根已知。这个技巧在圆锥曲线题中极为常用，能省下大量计算。
\end{gktip}

\par\medskip

\end{answer}
\begin{answer}
\gksolhead{22}{12 分}

\textbf{题目}　已知 $f(x)=\mathrm e^{x}-ax$ 和 $g(x)=ax-\ln x$ 有相同的最小值。(1) 求 $a$；(2) 证明：存在直线 $y=b$，与两曲线共有三个不同交点，且从左到右三个交点的横坐标成等差数列。

\par\medskip

\textbf{(1) 解}

\textbf{第一步，分别求两个最小值。}

\textbf{对 $f(x)=\mathrm e^{x}-ax$（定义域 $\mathbf R$）}：
\[f'(x)=\mathrm e^{x}-a\]

\begin{itemize}
\item 若 $a\le0$：$f'(x)>0$ 恒成立，$f$ 单调递增，\textbf{无最小值}（$x\to-\infty$ 时 $f\to-\infty$ 或 $0^{+}$，取不到最小）。
\item 若 $a>0$：由 $f'(x)=0$ 得 $x=\ln a$，且 $f$ 先减后增，最小值为
\end{itemize}
\[f(\ln a)=\mathrm e^{\ln a}-a\ln a=a-a\ln a\]

\textbf{对 $g(x)=ax-\ln x$（定义域 $(0,+\infty)$）}：
\[g'(x)=a-\frac1x\]

同理需 $a>0$，由 $g'(x)=0$ 得 $x=\dfrac1a$，最小值为
\[g\left(\frac1a\right)=a\cdot\frac1a-\ln\frac1a=1+\ln a\]

\textbf{第二步，令两最小值相等。}
\[a-a\ln a=1+\ln a\]

\textbf{第三步，解方程并证明唯一性。}

\textbf{观察}：$a=1$ 时左边 $=1-0=1$，右边 $=1+0=1$  是一个解。

\textbf{唯一性}：令
\[h(a)=a-a\ln a-1-\ln a,\qquad a>0\]

\[h'(a)=1-\left(\ln a+1\right)-\frac1a=-\ln a-\frac1a=-\left(\ln a+\frac1a\right)\]

令 $\phi(a)=\ln a+\dfrac1a$，则
\[\phi'(a)=\frac1a-\frac1{a^{2}}=\frac{a-1}{a^{2}}\]

$\phi$ 在 $a=1$ 处取最小值 $\phi(1)=0+1=1>0$，故 $\phi(a)>0$ 恒成立，从而
\[h'(a)<0\quad\text{对一切 }a>0\]

$h$ \textbf{严格递减}，故 $h(a)=0$ 至多一解，结合 $h(1)=0$ 知
\[a=1\]

\gkanswer{答案：$a=1$}

\par\medskip

\textbf{(2) 证明}

\textbf{第一步，发现两条曲线的内在联系（全题的钥匙）。}

$a=1$ 时：
\[f(x)=\mathrm e^{x}-x,\qquad g(x)=x-\ln x\]

在 $g$ 中令 $x=\mathrm e^{u}$（即 $u=\ln x$）：
\[g\left(\mathrm e^{u}\right)=\mathrm e^{u}-u=f(u)\]

即
\[\boxed{g(x)=f(\ln x)}\]

\textbf{推论}：$g(x)=b$ 的解与 $f(u)=b$ 的解通过 $x=\mathrm e^{u}$ \textbf{一一对应}。

\textbf{第二步，分析直线 $y=b$ 与两曲线的交点。}

由 (1)，$f_{\min}=f(0)=1$，$g_{\min}=g(1)=1$。

取 $b>1$，则：
\begin{itemize}
\item $f(x)=b$ 有两个解，记为 $u_1<0<u_2$（因 $f$ 在 $x=0$ 处取最小值 $1$）
\item $g(x)=b$ 的两个解为 $\mathrm e^{u_1}$、$\mathrm e^{u_2}$，满足 $0<\mathrm e^{u_1}<1<\mathrm e^{u_2}$
\end{itemize}

\textbf{四个横坐标}：$u_1$、$u_2$、$\mathrm e^{u_1}$、$\mathrm e^{u_2}$。

按大小关系定位：
\[u_1<0<\mathrm e^{u_1}<1<\mathrm e^{u_2},\qquad u_2>0\]

要使交点恰好为\textbf{三个}，必须有两个重合。观察可能性：
\begin{itemize}
\item $u_1$ 是唯一的负数，不可能与其它三个（都为正）相等；
\item $\mathrm e^{u_1}<1<\mathrm e^{u_2}$，二者不等；
\item 唯一可能：$u_2=\mathrm e^{u_1}$
\end{itemize}

\textbf{第三步，证明"三个交点"与"等差"是同一个条件。}

设 $u_2=\mathrm e^{u_1}$。由 $f\left(u_1\right)=f\left(u_2\right)$：
\[\mathrm e^{u_1}-u_1=\mathrm e^{u_2}-u_2\]

把 $\mathrm e^{u_1}=u_2$ 代入左边：
\[u_2-u_1=\mathrm e^{u_2}-u_2\ \Longrightarrow\ \mathrm e^{u_2}=2u_2-u_1\qquad(\ast)\]

此时三个交点的横坐标从左到右为
\[u_1<u_2\left(=\mathrm e^{u_1}\right)<\mathrm e^{u_2}\]

它们成等差数列的充要条件是
\[2u_2=u_1+\mathrm e^{u_2}\iff \mathrm e^{u_2}=2u_2-u_1\]

\textbf{这恰好就是 $(\ast)$！} 也就是说，只要三个交点的情形出现，等差数列\textbf{自动成立}。

\textbf{第四步，证明这样的 $b$ 确实存在。}

设 $u_1=t<0$，则 $b=f(t)=\mathrm e^{t}-t$。需要 $u_2=\mathrm e^{t}$ 也是 $f(x)=b$ 的根，即
\[f\left(\mathrm e^{t}\right)=b\ \Longrightarrow\ \mathrm e^{\mathrm e^{t}}-\mathrm e^{t}=\mathrm e^{t}-t\]

\[\mathrm e^{\mathrm e^{t}}-2\mathrm e^{t}+t=0\]

令
\[\psi(t)=\mathrm e^{\mathrm e^{t}}-2\mathrm e^{t}+t\]

计算两个端点：
\[\psi(0)=\mathrm e^{1}-2+0=\mathrm e-2\approx0.718>0\]
\[\psi(-1)=\mathrm e^{1/\mathrm e}-\frac2{\mathrm e}-1\approx1.4447-0.7358-1=-0.291<0\]

$\psi$ 在 $[-1,0]$ 上连续，由\textbf{零点存在定理}，存在 $t_0\in(-1,0)$ 使 $\psi\left(t_0\right)=0$。

取 $b=\mathrm e^{t_0}-t_0$，则直线 $y=b$ 与两曲线恰有三个不同交点，且横坐标成等差数列。$\blacksquare$

\begin{gktip}
\textbf{方法总结}：本题第 (2) 问是近年压轴题中构思最巧的一道，破题链条为：
1. \textbf{发现 $g(x)=f(\ln x)$}——这个恒等式把两条看似无关的曲线绑定，是全题的支点（提示来自 $a=1$ 时两函数的对称结构）；
2. 由此得出四个交点的横坐标为 $u_1,u_2,\mathrm e^{u_1},\mathrm e^{u_2}$，\textbf{用大小关系排除到只剩 $u_2=\mathrm e^{u_1}$ 一种重合可能}；
3. 惊喜地发现"三交点条件"与"等差条件"是\textbf{同一个等式}，两个要求一次满足；
4. 最后用零点存在定理保证存在性（不必解出 $t_0$ 的具体值——存在性命题只需证明存在）。

\textbf{考试策略}：第 (1) 问是常规的双函数求最值（可拿满分），第 (2) 问属于最高难度，建议及格线至 120 分段的学生只做 (1)。
\end{gktip}


\end{answer}